\section{Conclusion}
\label{sec:conclusion}
% Your version: 07_conclusion_draft.tex (future work and recommendations kept).

This project asked how the coupled dynamics of several interleaved buck-boost power shufflers in one string can be modelled and controlled. The answer is the $N$-module averaged model. It splits each shuffler into a shuffled mode and a circulating mode, and it gives the string's DC gain in closed form, $\vd \approx -\Vbar\uc$ and $\ic = (\Vbar/R_c)\ud$, with neighbour coupling below one per cent. Inverting that gain gives a decoupling integral controller on $[\vd, \ic]$: it commands the shuffled currents so the string equalises, and a partially shaded string reaches its maximum power point when the module voltages are equal, which the loop enforces.

The controller was carried from the analytical model through Simulink and QSPICE to the firmware, as one C core and one results pipeline, and validated on the bench and in the sun. On hardware it holds the module-voltage differences within $\pm0.010\V$ and the circulating currents within $\pm0.07\A$, limited by sensing noise. Its bandwidth follows the design $\tau = 1/(k\,2\pi\,150)$ to within the spread of the three simulations (0.4--1.2$\times$). Across 21 test windows on two days the loop closed from $k=0.1$ to 0.8 in the sun without a fault. The only instability was deliberately provoked ($k = 0.6$ at 300\,kHz), and the only trips came from live gain changes. With the top panel shaded to almost zero current, the shufflers delivered up to 99\W{} and 10.6\Wh{} over the afternoon to a 61\,V load that a plain series string could not have supplied. The averaged model's one real error, a fivefold over-estimate of the current-loop gain, was traced to a switching resistance $R_{sw} = 64\,C_p\fsw^3(L+M)^2$ and caught in QSPICE before any hardware ran. The result is a complete, reproducible system and a starting point for further work on MIMO control of DPP strings.

\textbf{Future work.} Scalability comes first: determine how the relationships change when more shufflers act in one string, or when a shuffler handles a varying number of panels, which needs boards that can be chained and driven independently. The gain calculation for a new switching frequency can now be streamlined, since $R_{sw}(\fsw)$ is a formula. Other controllers should be generated and compared against this one (LQR, MPC, an observer-based design, and the extended-Kalman-filter approach of \cite{kabalci_highly_2022}), as should a cost-function or machine-learning search for optimal gains. The firmware should move from its superloop to a real-time operating system, as FreeRTOS was planned, so that the control, sensing, safety and telemetry tasks are scheduled with explicit priorities and deadlines. Three firmware items were found in testing: refuse \code{C,APPLY} while the loop is running, add the COMP/DAC hardware over-current trip, and rework U10's isolator. Further work should also test other panel arrangements and dynamic configurations, and applications outside solar (mining, computing, industrial power and batteries), including whether the technology would win a solar car race.

\textbf{Recommendations.} The most immediate work is more data: stress-test the system on a very sunny day with every panel connected. To measure efficiency properly, add an MPPT algorithm, run it on the Raspberry Pi, which already controls the load over USB, and keep the load's limits in mind if the same set-up is reused. The electrical safety limits restrict which experiments can be run, so they should be designed around from the start. The system favours high switching frequencies, so GaN devices are worth revisiting. Mechanically, downsize the heatsink, improve the fasteners, reconfigure the PCB and redesign the enclosure, and consider consolidating the microcontroller and the microcomputer onto one board. Finally, hands-on development was the fastest route to understanding: anyone continuing the project should start from the GitHub repository, its manufacturing notes and its user guides.
